\thispagestyle{empty}
\begin{center}
{\Large\bf Lectures on}\\[5pt]
{\Large\bf  Forms of Higher Degree}
\vskip 1cm

{\bf  By}
\medskip

{\large\bf  J.I. Igusa}
\vfill


{\bf  Notes by}
\medskip

{\large\bf  S. Raghavan}
\vfill

  Published for the 

{\bf  Tata Institute of Fundamental Research,}

{\bf  Bombay}

{\bf  1978}
\end{center}


\eject

\thispagestyle{empty}

\begin{center}
{\Large\bf Lectures on}\\[5pt]
{\Large\bf  Forms of Higher Degree}
\vskip 1cm

{\bf  By}
\medskip

{\large\bf  J.I. Igusa}
\medskip


\vfill
{\bf  Notes by}
\medskip

{\large\bf  S. Raghavan}
\vfill

Published for the 

{\bf  Tata Institute of Fundamental Research,}

{\bf  Bombay}


\begin{figure}[H]
\centering{\includegraphics{logo.eps}}
\end{figure}

{\bf Narosa Publishing House}

New Delhi

1978 
\end{center}

\eject

\thispagestyle{empty}
\begin{center}
  \textbf{Author}\\[5pt]
  {\large\bf Jun-Ichi Igusa}\\
  Department of Mathematics\\
  Johns Hopkins University\\
  Baltimore, Maryland 21218\\
  U.S.A.\\
  \vfill

  \copyright\, {\bf Tata Institute of Fundamental Research}

  {\bf 1978}
  \vfill


\parbox{0.7\textwidth}{No part of this book may be reproduced in 
any form by print, microfilm or any other means without written 
permission form the Tata Institute of Fundamental Research, 
Colaba, Bombay 400 005.}
\vfill

  Printed by
 
  {\bf K.P. Puthran} at the Tata Press Limited,

  414 Veer Savarkar Marg, Bombay 400 025 and 
\medskip

  Published by
 
 {\bf  N.K. Mehra Narosa Publishing House,}
 
  2/35 Ansari Road, Daryaganj, New Delhi-110 002.

  \textbf{PRINTED IN INDIA}
\end{center}

\eject


\thispagestyle{empty}

\chapter{Foreward}


THE FOLLOWING lecture notes are based on my lectures at the Tata
Institute of Fundamental Research from the middle of January to early
March of 1978. My objective was not to give a systematic exposition of
the theory of forms of higher degree but only to introduce young
mathematicians to a new approach to this area. It is my view that
although this theory is still at an early stage, it has all the
features of becoming an important branch of Mathematics. If the reader
finds these lecture notes intelligible, then it is largely due to the
painstaking efforts of Professor Raghavan. I would like to thank him
for his excellent collaboration to bring the notes into this final
form.
\vskip 2cm

\noindent
9 March 1978 \hfill {\large\bf Jun-Ichi Igusa}
